% Generated draft: requires mathematical and source review.
\documentclass[11pt]{article}
\usepackage[margin=25mm]{geometry}
\usepackage{fontspec}
\setmainfont{Latin Modern Roman}
\usepackage{amsmath,amssymb,mathtools,unicode-math}
\setmathfont{Latin Modern Math}
\usepackage{enumitem,needspace,xurl,hyperref}
\setlength{\parindent}{0pt}
\setlength{\parskip}{5pt}
\setlength{\emergencystretch}{3em}
\newcommand{\N}{\mathbb{N}}
\newcommand{\Z}{\mathbb{Z}}
\newcommand{\Q}{\mathbb{Q}}
\newcommand{\R}{\mathbb{R}}
\newcommand{\C}{\mathbb{C}}
\newcommand{\F}{\mathbb{F}}
\newcommand{\A}{\mathbb{A}}
\newcommand{\PP}{\mathbb{P}}
\newcommand{\E}{\mathbb{E}}
\newcommand{\eps}{\varepsilon}
\newcommand{\dd}{\,\mathrm{d}}
\begin{document}
\section*{UnsolvedMath ID: 30001937}
\textbf{Source problem code:} catalogue entry
\subsection{Definitions and mathematical statement}
A convex cone in \(\R^m\) is a convex set closed under multiplication by nonnegative real scalars. It is full-dimensional if its interior in \(\R^m\) is nonempty. Given a cone \(K\subseteq\R^m\) and a set \(C\subseteq\R^n\), a \(K\)-lift of \(C\) consists of an affine subspace \(L\subseteq\R^m\) and a linear map \(\pi:\R^m\to\R^n\) such that
\[
C=\pi(K\cap L).
\]
The affine subspace is not required to pass through the origin. Equality here means exact equality of sets; no closure operation or approximation is included.

The original formulation takes \(K\) to be a full-dimensional closed convex cone and \(C\) to be an arbitrary convex set, and asks when such a lift exists. The supplied corrected formulation focuses on the minimum-size question and calls \(C\) a cone. The supplied verification note asserts that the report proves the preceding general characterization and leaves the minimum-size question open; this assertion has not been independently verified here.

For a specified nested family \((K_k)_{k\in I}\), with each \(K_k\subseteq\R^{m_k}\), and a specified target cone \(C\subseteq\R^n\), define the set of feasible indices by
\[
\mathcal F(C;(K_k))=
\left\{k\in I:\begin{array}{l}
\text{there exist an affine subspace }L\subseteq\R^{m_k}\\
\text{and a linear map }\pi:\R^{m_k}\to\R^n\\
\text{such that }C=\pi(K_k\cap L)
\end{array}\right\}.
\]
The minimum-size question is to determine, in terms of the specified family and target,
\[
\min\mathcal F(C;(K_k)),
\]
whenever this minimum exists, and to determine whether a feasible index exists at all. If the feasible set is empty, no member of the family provides such a lift. If the index order does not ensure existence of a least feasible index, that issue must also be resolved.

\textit{Unresolved scope.} The record supplies neither a particular family nor a particular target cone. It does not specify the index set and its order, the meaning of nesting when ambient dimensions vary, or whether every \(K_k\) must satisfy the full-dimensionality, closedness, and convexity hypotheses imposed on \(K\) in the original formulation. It also does not explicitly specify convexity, closedness, or full-dimensionality of the target cone in the corrected formulation. These qualifications must be settled before this general problem schema becomes a fully specified minimum-size problem.
\subsection{Short English statement}
For a chosen family of cones and a chosen target cone, find the least index for which an affine slice of the indexed cone maps linearly onto the target. The source leaves the family, target, and several conventions unspecified. Its original formulation allows arbitrary convex target sets, whereas its corrected formulation asks about target cones. This is a statement-only draft requiring human mathematical review.
\subsection{Sources}
\begin{itemize}
\item Supplied UnsolvedMath local registry record, permanent ID \(30001937\), title \emph{Minimum Size of Cone Lifts}, collection \emph{Oberwolfach Reports Open Problems}, exact source problem code \texttt{catalogue entry}. No individual proposer is supplied; the proposed year is recorded as 2011.
\item Supplied report citation: \emph{Combinatorial Optimization (2012)}, Oberwolfach Reports, pp.\ 3003--3076, DOI \texttt{10.4171/owr/2011/53}. Supplied source URL and literature reference labelled ``original OWR report'': \url{https://doi.org/10.4171/owr/2011/53}. The background attributes the problem to mathematical workshop discussions documented in Oberwolfach Reports. The report has not been consulted here.
\item Supplied literature assessment, checked on exactly \texttt{2026-08-21}: the record labels the problem open and reports that no general classification or resolution was found for the cone-lift representation and least-size questions. This is supplied dated curation, not an independently verified literature assessment. The background explicitly requests independent verification of time-sensitive status and cited claims.
\item The supplied statement status is \texttt{corrected\_verified}. Its verification note asserts that the official report proves the preceding general characterization of cone lifts and leaves only the minimum-size question open. This supplied assertion requires reconciliation with the broader literature assessment; the registry label does not itself establish mathematical correctness.
\end{itemize}
\end{document}
